\section{Purpose, principal results and conventions}
\label{intro:section}

The aim is to add proved identities to the current research programme,
while preserving the distinctions between an ordinary integral, a specified
finite part, a formal reconstruction theorem and a numerical observation.
The starting points are the manuscript and incoming directories named by
the project maintainer \cite{RepoManuscript,RepoIncoming}. All twelve ZIP
archives present in the retrieved incoming listing were downloaded and
their mathematical source files inspected by topic. Ninety-two manuscript
chapter files were also retrieved. This is a targeted proof and status
audit of the directions used here, not a claim to have checked every line
of every chapter. File hashes and the retrieval scope are supplied in
\file{data/source_manifest.json}.

The newest source matters. In particular, the report
\emph{Polygamma Moments and Stieltjes Resonances} already completes gap
eight, treats all positive monomial scales for two Stieltjes factors,
and evaluates a substantial joint-polygamma range
\cite{PolygammaResonances}. The expanded incoming version of
\emph{Universal Gamma Normal Forms} already gives the coefficient of
every canonical monomial \cite{UniversalIncoming}. Those results are
inputs here. The new work starts with the next fixed gap, additional
collision factors, a third mass cluster, and the still missing
nonrecursive primitive and infinite-jet convergence statements.

\subsection{A map of the contributions}

\begin{longtable}{@{}>{\raggedright\arraybackslash}p{0.21\textwidth}>{\raggedright\arraybackslash}p{0.38\textwidth}>{\raggedright\arraybackslash}p{0.34\textwidth}@{}}
\toprule
Direction & Result proved here & Precise remaining limit\\
\midrule
\endhead
Polygamma collisions & A finite Hermite-divided-difference correction
for any number of factors and arbitrary positive shifts;
Theorem~\ref{pc:uniform}. & The general merged product is retained as a
specified finite part. Positive Stieltjes indices require further work.\\[5pt]
Evaluated constants & The complete one-trigamma, centered-digamma family,
Theorem~\ref{pc:all-powers}; nonresonant and resonant triples. & The scale
hypotheses determine which constant is selected.\\[5pt]
Multiple mass clusters & A classical-polylogarithm divided difference
for every labelled mixed derivative, and a depth-three primitive for
$B^r,C,D$; Theorems~\ref{tri:thm:mixed} and~\ref{tri:thm:three}.
& The Gamma result is an ordinary spectral sum; no general finite
arithmetic closure is proved.\\[5pt]
Harmonic reconstruction & Every unequal even pair with gap ten has the
optimal formal prefix; Theorem~\ref{gt:main}. & The unrestricted
unequal-even-pair conjecture remains open.\\[5pt]
Inverse coefficients & An $(m+1)$-square bordered determinant at fixed
half-gap $m$; Theorem~\ref{gt:inverse}. & Its entries still use a finite
causal recurrence.\\[5pt]
Gamma--Stieltjes primitives & A nonrecursive finite formula for every
compensator $H_m$ and locally convergent normal and primitive series;
Theorems~\ref{hp:closed}, \ref{hp:analyticN} and~\ref{hp:analyticH}.
& Sharp convergence domains, coefficient signs and arithmetic
denominators remain separate questions.\\
\bottomrule
\end{longtable}

The contact and Gamma-primitive arguments are analytic or algebraic
proofs at every order. The gap-ten theorem has a finite exact
computational component, fully exposed below. Numerical checks of
polylogarithms and shifted integrals are diagnostics of the displayed
formulas; they are not substitutes for the proofs.

\subsection{Two representative identities}

The most compact collision identity can be stated before introducing
the full local correction. Let $\alpha_1>\alpha_2>\alpha_3>0$ be
linearly independent over $\mathbb Q$, and let $r\in\{1,2,3\}$.
Then
\begin{equation}\label{intro:triple}
\boxed{\displaystyle
\CT_{\epsilon\downarrow0}\int_0^1
 \psi'(x+\epsilon^{\alpha_r})
 \prod_{j\ne r}\psi(x+\epsilon^{\alpha_j})\,dx
 =\begin{cases}
  \zeta(3),&r=1,\\
  0,&r=2,\\
  -\zeta(3),&r=3.
 \end{cases}}
\end{equation}
Here the smallest shift carries the largest exponent. Every integral
is ordinary for $\epsilon>0$; the operation outside it is a specified
asymptotic constant. The fully evaluated subtraction formula in
Section~\ref{pc:section} also expresses these constants as ordinary
renormalized limits. No fitting of divergent coefficients is needed.

For the higher family put $\Psi=\psi+\gamma$ and
$f(x)=x(\psi(1+x)+\gamma)=\sum_{k\ge2}(-1)^k\zeta(k)x^k$.
With the analogous nonresonant ordering of $N$ scales, the complete
constant with the derivative in position $r$ is
\begin{equation}\label{intro:allN}
\boxed{\displaystyle
T_{N,r}=-[x^N]\left\{\frac{(f-1)^N}{N}
                     +f^{N-r}(f-1)^{r-1}\right\}.}
\end{equation}
Thus a single coefficient formula yields a weight-$N$ zeta polynomial
for every $N$ and every derivative position. Its adjacent differences
and its zero row sum are identities, rather than patterns inferred from
the first rows.

The representative spectral identity has a different character. For
the explicitly defined ordinary simplex face $I_k$ and Gamma field $G_a$
of Section~\ref{tri:sec:main}, all $r,m\ge1$ satisfy
\begin{equation}\label{intro:mixed}
\partial_{C_1}\cdots\partial_{C_m}
 I_{r+m}(G_a;B^r,C_1,\ldots,C_m)
 =(-1)^r r!\,\mathfrak g_{r,a,B}[C_1,\ldots,C_m],
\end{equation}
where
\[
\mathfrak g_{r,a,B}(C)=\sum_{n\ge0}
 \frac{\operatorname{Li}_r((C-B)/((n+a)^2+C))}{(n+a)^2+C}.
\]
The number of independent masses does not increase the depth of the
polylogarithm in the differentiated formula. The finite three-cluster
primitive reconstructs the information on the coordinate edges as well
as the common mixed derivative.

\subsection{Conventions used throughout}

We write $\psi^{(r)}$ for the $r$th argument derivative of the digamma
function, with $\psi^{(0)}=\psi$. The Hurwitz function is
$\zeta(s,a)$, and a superscript on $\zeta$ denotes an order derivative
unless an argument derivative is explicitly displayed. The Stieltjes
normalization is
\begin{equation}\label{intro:Stieltjes}
\zeta(1+s,a)=\frac1s+
    \sum_{q\ge0}\frac{(-1)^q}{q!}\gamma_q(a)s^q,
\qquad \gamma_0(a)=-\psi(a).
\end{equation}
These conventions agree with the standard Gamma and Hurwitz-zeta
identities \cite{DLMF,DLMFHurwitz}.
Generalized harmonic numbers are
$H_n^{(s)}=\sum_{j=1}^n j^{-s}$, with $H_0^{(s)}=0$.
The Bernoulli convention is $B_1=-1/2$.

For a locally finite power--logarithm expansion,
$\CT_\epsilon$ selects the coefficient of
$\epsilon^0(\log\epsilon)^0$. The merged finite part at the lower
endpoint is always taken in the displayed coordinate $x$:
$\CT_{\delta\downarrow0}\int_\delta^1$.
It is not invariant under an arbitrary change of cutoff coordinate.
All endpoint subtractions and all shifts are therefore kept explicit.

For divided differences, one node gives evaluation and repeated nodes
give derivatives divided by the corresponding factorial. Cumulative
colors for multiple polylogarithms are defined in
\eqref{tri:eq:cumulative}; this convention makes the relevant real
branches and zero-color limits transparent. Mass derivatives are taken
in independently labelled variables before imposing collisions.

Finally, formal rank in the harmonic reconstruction problem means
rational matrix rank on declared symbols. It does not imply arithmetic
independence of the numerical zeta or moment values assigned to them.
The new conclusions are established relative to the inspected project
sources. Classical homotopy, divided-difference and finite-part methods
are credited; a bounded literature search cannot certify worldwide
priority for the assembled identities.
